\paragraph{Las dos trazas del término mixto.}
Pongamos \(D=H_*D_*-S_*\). La ciclicidad de la traza da la identidad en el padre
\[
 Q_c=\tfrac12\sum_{u,v}(S_*)_{uv}(m^2D)_{vu}
       +\tfrac1{4m^2}\sum_{u,v}(m^2D)_{uv}(m^2D)_{vu}.
\]
Cada casilla de \(F\setminus\{c\}\) tiene tres líneas retenidas, que dan seis pares
ordenados en la primera suma. En la segunda se integran dos líneas libres.
Para expresar \(m^2D\) en el perfil escribamos \(\mathcal K(v,w)=E(v,w)+\Pi_{vw}\),
\(\mathcal V(v,j)=\pesodirecto(v\cap j)+\Pi_{vj}\), \(j\in \familias(c)\). Para líneas retenidas \(v,u\),
sea \(\mathcal Y(v,u)=-\pesodirecto(z)T_z\) si se cortan en \(z\in E'\), y cero en otro caso.
En esta fórmula \(T_z\) usa los mismos indicadores que \(\pesodirecto(z)\).
Finalmente, para \(w\ne u\) retenidas, sea \(\mathcal S(w,u)=-\pesodirecto(i')\) cuando
\(w\cap u=i'\in F\setminus\{c\}\), y cero en los demás pares. Definamos
\begin{align*}
 \mathcal B(v,u)&=\mathcal Y(v,u)+\Pi_{vu}
                    +\sum_{j\in \familias(c)}\mathcal V(v,j)\mathcal V(u,j),\\
 \mathcal E_2(v,u)&=\mathcal B(v,u)+\sum_w\mathcal K(v,w)\mathcal S(w,u)
                     +\int\mathcal K(v,w)\mathcal B(w,u)\,d\rho(w).
\end{align*}
Estas tres piezas proceden, con \(B_*=D_*-S_*\) y \(K_*=H_*-I\), de
\(m^2D=m^2B_*+(mK_*)(mS_*)+m^{-1}(mK_*)(m^2B_*)\).
En la suma discreta se toma cada \(i'=u\cap j\), \(j\in \familias(c)\), y cada una de
sus dos líneas transversales \(w\). Su disponibilidad está en \(\pesodirecto(i')\).
En la integral sí se marca \(w\), pues su disponibilidad ya está en \(d\rho(w)\).

Escribamos \(R(i)=\{(u,v):u,v\in \familias(i)\setminus\{\lambda\},\ u\ne v\}\) para una
casilla de la fibra \(\lambda\). Aplicando \(\psi\) a los monomios se obtienen
\begin{align*}
 L_1&=-\tfrac12\left\langle\int_F
          \sum_{(u,v)\in R(i)}\mathbb E\mathcal E_2(v,u)\right\rangle_c,\\
 L_2^{\rm int}&=\tfrac14\left\langle\iint
 \mathbb E\big[\mathcal E_2(u,v)\mathcal E_2(v,u)-\mathcal Y(u,v)\mathcal Y(v,u)\big]
                 \,d\rho(u)\,d\rho(v)\right\rangle_c .
\end{align*}
Las dos copias de una línea integrada internamente tienen variables libres separadas.
El bloque \(\mathcal Y\mathcal Y\) se sustrae antes de integrar. En el tablero cada
\(i\in E'\) aporta doce pares ordenados al bloque \(YY\), de modo que su contribución
es \(+3m^2\sum_{E'}y_i^2\), exactamente la opuesta a la de \(P_c\).
También se puede comprobar en el continuo. Para cada par de familias, la medida de los
cruces válidos por su disponibilidad es \(U_DU_A\,di\), con las hojas y el jacobiano del
cruce incluidos. Multiplicando por \(W(i)^2\) resulta \(\nu_1(i)W(i)\).

\paragraph{Cómo extraer \(\Delta_1\) y \(\Delta_2\).}
Las filas tabuladas son \(Q_1=L_1-\Delta_1\) y \(Q_2=L_2^{\rm int}-\Delta_2\).
Damos las correcciones por separado para poder reconstruir cada entrada. En todas las
sumas siguientes un cruce debe existir en el cuadrado y en la clase de paridad indicada.
Una igualdad de líneas se exige como identidad en los parámetros libres de esa clase.
Las coincidencias accidentales forman conjuntos de medida nula.

Para \(\Delta_1\), fijemos \(c,i,\lambda,(u,v)\in R(i)\) como en \(L_1\), y una
línea libre \(w\) integrada con \(d\rho(w)\). Pongamos \(z=v\cap w,\ p=w\cap u\).
Para cada \(a\in \familias(p)\setminus\{w,u\}\) y \(j\in \familias(c)\), pongamos \(t=a\cap j\).
Se marcan \(\familias(c)\cup \familias(i)\cup\{w\}\).
El monomio directo de \(\mathcal K(v,w)\mathcal Y(w,u)\) aporta la diferencia
\[
 \Delta_1=\tfrac12\left\langle\int_F\sum_{(u,v)\in R(i)}
 \int\sum_{a,j} W(z)W(p)W(t)
 \big[\psi(\{z,p,t\})-\psi(\{z\})\psi(\{p,t\})\big]\,d\rho(w)\right\rangle_c .
\]
Se toman los tres cruces distintos, con \(p\in E'\). La diferencia cuenta la transversal
común de \(z,t\) sólo una vez. Si es \(\ell\), el corchete vale
\(\psi(\{z\})\psi(\{p,t\})(U_\ell^{-1}-1)\).
Las tres partes directas llevan signo negativo en el producto de la primera traza
antes de aplicar su prefactor \(-1/2\).
Por ejemplo, en coordenadas relativas, para \(c=(0,0)\), \(i=(0,s)\), \(u=C_s\),
\(v=D_{-s}\), \(w=R_h\), \(a=A_{h+s}\), \(j=R_0\), resultan \(z=(h,h+s)\),
\(p=(h,s)\), \(t=(0,h+s)\). La columna \(h+s\) aparece en \(z\) y \(t\).

Para \(\Delta_2\), marquemos \(\familias(c)\cup\{u,v\}\), con \(u,v\) libres. Cuando existe el cruce,
el peso de sus transversales es
\[
 K_{ab}=W(a\cap b)\prod_{\ell\in \familias(a\cap b)\setminus\{a,b\}}U_\ell .
\]
Para cada \(j\in \familias(c)\) tomamos \(p_2=v\cap j\),
\(w_2\in \familias(p_2)\setminus\{v,j\}\), \(t_2=u\cap w_2\), \(k_2=K_{vj}K_{uw_2}\).
Los dos patrones posibles para los otros dos vértices dan
\[
\begin{array}{c|c|c|c}
 &p_1,\ t_1 &k_1&\text{multiplicidad }s\\ \hline
 \mathcal K\mathcal S\ \text{por}\ \mathcal K\mathcal S&
 u\cap j,\ v\cap w_1 & K_{uj}K_{vw_1}&1\\
 \mathcal Y\ \text{por}\ \mathcal K\mathcal S&
 v\cap u,\ a\cap j & K_{vu}K_{aj}&2
\end{array}
\]
En la primera fila \(w_1\) recorre \(\familias(p_1)\setminus\{u,j\}\). En la segunda \(a\)
recorre \(\familias(p_1)\setminus\{v,u\}\). En cada patrón se conservan cuatro vértices distintos.
Para cada pareja de cruces que define \(k_d\), \(d=1,2\), reunamos las transversales de
sus dos cruces, es decir, las líneas cuyos factores \(U\) aparecen en \(k_d\).
Si los dos conjuntos comparten una línea \(\ell\) fuera de las marcadas, se incluye el
sumando siguiente. Los cuatro cruces forman entonces un cuadrado y la transversal es única.
\[
 \Delta_2=\tfrac14\left\langle\iint\sum_{\text{patrones admitidos}}
          s\,k_1k_2(U_\ell^{-1}-1)\,d\rho(u)\,d\rho(v)\right\rangle_c .
\]
Se puede obtener la lista desarrollando los productos de \(\mathcal E_2\).
Un peso \(\pesodirecto\) aporta las cuatro coordenadas de su casilla menos las líneas ya marcadas.
Para comparar dos etiquetas de una misma familia se restan sus expresiones afines en
las coordenadas libres. Sólo una diferencia idénticamente nula aporta el factor
\(U_\ell^{-1}-1\). Para hacer la comprobación basta comparar el término constante y
los coeficientes de cada parámetro en las cuatro etiquetas \(r,q,r-q,r+q\). En
\(\Delta_1\) se obtienen 24 elecciones de familias y 80 al distinguir las clases admisibles.
En \(\Delta_2\) quedan 24 elecciones en cada fila de la tabla. Por ejemplo, en la primera,
con \(j=R_0,u=A_s,v=D_t,w_1=C_s,w_2=C_{-t}\), los dos vértices exteriores
\((s+t,s),(s+t,-t)\) comparten la fila \(s+t\). En la parte mixta, intercambiar las copias da el mismo
integrando con \(u,v\) permutadas, de donde procede el \(2\) de la tabla.
Finalmente,
\begin{align*}
 d(x)={}&G_{\rm cof}+f2E+f2z+2f2F+f3F+f3c+f4+f5c+f5F\\
        &+f6+f7+C_{3,\rm cuad}+Q_1+\Delta_1+Q_2+\Delta_2 .
\end{align*}
